\chapter*{概要}
\thispagestyle{empty}
\enlargethispage{\baselineskip}

\begin{center}
{\large 視覚情報に基づく自律走行モビリティの\\ラップタイムを決定する要因の分析}
\end{center}

{\small\renewcommand{\baselinestretch}{1.1}\selectfont
カメラと内界センサだけで走行する小型モビリティは，LiDARと地図に依存せず，地図の整備が難しい構内や屋内外の境界で使える．AI-Formulaは，LiDARと自作地図の使用を禁止し，信号・交差点・障害物を含むコースの3周の時間を競う競技会であり，この条件を模している．実用には安全な範囲で速度を上げることが求められるが，従来の評価は走行の成功率や経路の推定精度が中心であり，ラップタイムがどの要素に由来し，速度を上げたときに何が短縮を妨げるかは示されていない．

本研究は，この競技に向けて開発した，カメラ画像からの経路生成，曲率に基づく速度計画，経路追従制御からなるシステムについて，ラップタイムを要素別の時間差に分け，上限速度を上げたときの短縮を妨げているものを明らかにする．方法は3つである．第一に，走行した軌跡の上で，車両で確認した性能だけを与えた基準時間から実測のラップタイムまで速度列を順に置き換え，上限速度，速度計画，経路生成，速度の追従，停止に起因する時間差に分ける内訳を定義する．第二に，これを上限速度4.0，5.0，6.0\,m/sの実車の周回走行9本に適用する．第三に，同定した車両応答と経路生成の揺らぎを組み込んだ閉ループシミュレーションで，設定値と要素を変えた条件を比べる．

実車では，上限速度を4.0から6.0\,m/sに上げると，2周目のラップタイムは赤コースで20.4\,s，青コースで14.7\,s短くなった．青コースでは，この短縮は上限速度による時間差の減少より6.6\,s小さく，そのうち6.1\,sは速度計画の横加速度の上限と分岐での速度抑制によるもので，経路生成と速度の追従に由来する時間差の増加は0.1〜0.2\,sであった．速度の追従の影響は，時間ではなく，速度計画の上限2.0\,m/s$^2$を超える横加速度（最大3.4\,m/s$^2$）として現れた．シミュレーションは，上限速度を変えたときのラップタイムの変化を実車と0.6\,s以内で再現し，青コースの6.0\,m/sで，横加速度の上限を3.0\,m/s$^2$に上げると2.75\,s，分岐での速度抑制を外すと3.10\,s短くなった（二つの短縮は足し合わせられない）．比較のために設けた，局所経路計画の合流距離を15\,mに固定する条件では，このシミュレーションの青コースの完走が上限速度5.0，6.0\,m/sで30試行中10に下がった．

以上から，上限速度6.0\,m/sまでの範囲で，短縮を減らしているのは速度計画の設定値であり，経路生成と速度の追従に由来する時間差の増加は0.1〜0.2\,sである．これらの影響は，時間差ではなく，横加速度の超過と分岐での経路の切替の遅れとして現れ，設定値はこれらに対する余裕として働いている．

\vspace{0.5\zh}
\noindent キーワード: 自律走行モビリティ，AI-Formula，ラップタイム，速度計画，経路追従，閉ループシミュレーション
}

\chapter*{Abstract}
\thispagestyle{empty}
\enlargethispage{\baselineskip}

\begin{center}
{\large Analysis of the Factors Determining the Lap Time\\of a Vision-Based Autonomous Vehicle}
\end{center}

{\small\renewcommand{\baselinestretch}{1.0}\selectfont
Small vehicles that drive only with a camera and proprioceptive sensors do not depend on LiDAR or prebuilt maps. AI-Formula, which prohibits LiDAR and self-made maps and ranks teams by the time of three laps with traffic signals, intersections, and obstacles, models this condition. For practical use the speed must be raised safely, but evaluations have focused on success rates and path accuracy, not on where the lap time comes from.

This thesis divides the lap time of a system developed for this competition, consisting of camera-based path generation, curvature-based velocity planning, and path following, into time differences by component, and identifies what prevents the lap time from decreasing when the upper speed limit is raised. First, on the driven trajectory, the velocity profile is replaced step by step from the one given only by the confirmed performance of the vehicle to the measured one, dividing the lap time into a reference time and time differences due to the upper speed limit, velocity planning, path generation, speed tracking, and stopping. Second, this is applied to nine real lap runs at 4.0, 5.0, and 6.0\,m/s. Third, a closed-loop simulation compares conditions with changed settings or components.

Raising the upper speed limit from 4.0 to 6.0\,m/s shortened the second-lap time by 20.4\,s on the red course and 14.7\,s on the blue course. On the blue course, the shortening was 6.6\,s smaller than the decrease in the time difference due to the upper speed limit; 6.1\,s of this came from the lateral acceleration limit and the branch speed limit of the velocity planning, while path generation and speed tracking added 0.1 to 0.2\,s. Speed tracking appeared as lateral acceleration above the planning limit of 2.0\,m/s$^2$ (up to 3.4\,m/s$^2$), not as time. The simulation reproduced the change of the lap time with the upper speed limit within 0.6\,s; on the blue course at 6.0\,m/s, a lateral acceleration limit of 3.0\,m/s$^2$ shortened the lap by 2.75\,s and removing the branch speed limit by 3.10\,s (not additive). In a condition made for comparison, fixing the merge distance to 15\,m, simulated completions on the blue course fell to 10 of 30 at 5.0 and 6.0\,m/s.

Up to 6.0\,m/s, the shortening is reduced by the velocity planning settings, while path generation and speed tracking added only 0.1 to 0.2\,s. Their effect appeared as lateral acceleration above the limit and the delay of the path switch at the branch, against which the settings act as margins.

\vspace{0.5\zh}
\noindent Keywords: Autonomous vehicle, AI-Formula, Lap time, Velocity planning, Path following, Closed-loop simulation
}
